\documentclass[11pt]{article}
\usepackage[utf8]{inputenc}
\usepackage[T1]{fontenc}
\usepackage{lmodern}
\usepackage[margin=1in]{geometry}
\usepackage{amsmath,amssymb}
\usepackage{textcomp}
\usepackage[dvipsnames]{xcolor}
\usepackage{enumitem}
\usepackage{parskip}
\usepackage{titlesec}
\usepackage[colorlinks=true,linkcolor=blue,citecolor=blue,urlcolor=blue]{hyperref}
\usepackage{natbib}

\titleformat{\section}{\large\bfseries}{\thesection}{1em}{}
\titleformat{\subsection}{\normalsize\bfseries}{\thesubsection}{1em}{}
\titleformat{\subsubsection}{\normalsize\bfseries\itshape}{}{0em}{}

\title{\textbf{The Changelog Is Not the Change:\\ Measuring Silent Behavioral Drift in Deployed LLM APIs}}
\author{}
\date{}

\newcommand{\ph}{\textit{[placeholder --- fill in after experiments are run]}}

\begin{document}
\maketitle
\vspace{-2em}

\begin{quote}
\small\textit{Assembled from the six confirmed stage outputs of an \texttt{auto-research-discussion} run (seed topic B3.3), post research-taste audit (\textcolor{ForestGreen}{\textbf{PASS}}). Status: \textbf{specification / paper draft --- no experiments have been run.} All numbers below are pre-registered design targets, not results.}
\end{quote}

\begin{abstract}
Vendors of deployed AI systems update their APIs on a schedule they alone control, and the changelog describing what changed is a document they alone author --- so the completeness of that record is unverifiable from inside it. We show that this is not a hypothetical gap: existing black-box audit tools can detect that a served model's behavior shifted --- the closest even reports suspected undisclosed changes at scale, by presuming that any change on a version-pinned endpoint is unannounced --- but none of them aligns a detected change against the vendor's own disclosure text to determine what was documented and what was not. We build a pipeline that runs a category-stratified battery of behavioral probes against paired pre/post-update API snapshots, separates genuine cross-version change from ordinary sampling noise using a calibrated same-version null and an effect-size floor, and aligns each detected change to the vendor's own changelog claims by semantic entailment, calibrated against human labels and a string-matching floor. Probe categories are anchored to independently published benchmarks rather than invented after the fact, so any resulting concentration pattern cannot be a taxonomy artifact; version boundaries are attributed using a data-driven fingerprint discontinuity rather than assumed from announcement dates; and the resulting silent-change rate is reported honestly as a coverage-conditioned lower bound, validated for stability rather than claimed as a point estimate. We evaluate this design across real vendor API transitions, comparing against a same-version control and a naive string-matching aligner, to report what fraction of measured behavioral change silently escapes vendor disclosure, and where it concentrates.
\end{abstract}

\section{Introduction}

Large language model APIs change on a schedule the vendor alone controls, and the record of what changed is a document the vendor alone authors. Integrators build products against a version they never chose to freeze, auditors evaluate systems whose behavior may already differ from what they last tested, and regulators drafting ``ongoing monitoring'' mandates for AI systems are, in practice, asking a party to monitor itself. The changelog is the one artifact everyone treats as the record of what changed between releases --- which makes it worth asking, directly, how complete that record actually is.

Yet a basic question about that arrangement remains unanswered: how much of what actually changed never appears in that record at all? A vendor who ships a behavioral change and discloses it, and a vendor who ships the same change and says nothing, both publish a changelog --- and nothing about the changelog's own text can tell an outside observer which vendor did which. This is not a minor documentation gap. If a substantial share of real behavioral change is silent, then any accountability process that treats the changelog as a complete record --- attributing an incident, a regression, or a safety-relevant shift to ``whatever the vendor said changed'' --- silently inherits a gap it has no way to see, let alone bound.

The nearest existing work approaches adjacent pieces of this question without closing it. The closest is Chen, Zaharia, and Zou's demonstration \citep{Chen2023ChatGPTBehavior} that GPT-3.5 and GPT-4 drifted substantially in behavior between two dated snapshots on a fixed task battery --- direct evidence that undisclosed drift occurs, but the study never cross-references OpenAI's own release notes to say which of the observed changes were documented and which were silent, which is precisely the missing step this paper supplies. Other black-box audit tools solve adjacent but distinct problems: continuous logprob-tracking methods \citep{LogProbTracking} detect that an endpoint's output distribution shifted --- and, in a months-long field deployment, have already flagged dozens of suspected undisclosed changes across production endpoints, ``undisclosed'' there being a presumption inherited from version pinning rather than a determination made by reading any disclosure text --- and rank-based uniformity tests \citep{RankUniformity} detect that a served model was quietly substituted for another; but both stop at the change signal itself, with no attempt to attribute the shift to a behavior category or align it against the changelog the vendor actually published. Fairness-specific audits \citep{BiasSimilarity} measure whether a tuned release departs from its base model, but only within a stereotype/bias construct, never against changelog text. And governance surveys \citep{TAIG} catalog the resulting monitoring gap as a policy problem without building or running the measurement themselves. Every ingredient for this paper's question exists somewhere in this landscape; no prior work puts them together against the vendor's own disclosure text.

Turning this gap into a measurable claim surfaces four distinct difficulties. Individual LLM outputs are stochastic even within a single fixed version, so a genuine cross-version change first has to be separated from ordinary sampling noise, load-dependent routing, and prompt sensitivity --- and the same-version noise floor is itself a moving target, not a fixed constant. Vendor changelogs are unstructured, variable-granularity prose (``improved reasoning,'' ``safety updates''), so deciding whether a detected behavioral shift is actually entailed by a changelog claim is an open semantic-matching problem, not a lookup. Any resulting claim about \emph{which behavior categories} concentrate silent change is only trustworthy if the category boundaries themselves were not chosen after seeing which probes happened to shift --- otherwise a ``concentration'' finding would be circular. And because API-served models are updated via rolling deploys, undated changes, and A/B routing rather than clean discrete releases, attributing a detected diff to the specific release window a changelog entry describes is itself an inference under uncertainty, not a given.

Our approach turns each of these difficulties into a specific, checkable design choice rather than an assumption. A calibrated diff detector requires both a corrected statistical test and a minimum effect size to call a probe changed, and its behavior is checked directly against a same-version control where no real change exists. A semantic-entailment aligner, not a keyword matcher, decides whether a changelog claim actually covers a detected diff, calibrated against human adjudication and benchmarked against how much a naive string-matching baseline could have gotten right on its own. Behavior categories are anchored to benchmarks that were designed and published for unrelated purposes, so a concentration finding cannot be an artifact of how the categories were drawn. Version boundaries are attributed from a continuously logged, data-driven fingerprint discontinuity rather than assumed from a vendor's announcement date. And because no independent ground truth of ``everything that actually changed'' in a release will ever exist, the resulting silent-change rate is reported explicitly as a lower bound and validated for stability under independent probe subsets, rather than presented as a number more precise than the evidence supports.

We evaluate this design across real version transitions from multiple LLM API vendors, comparing the cross-version pipeline against a same-version null control and a naive string-matching changelog aligner, and reporting the resulting documented-vs-silent rate per behavior category with bootstrap confidence intervals.

In summary, we contribute:
\begin{itemize}[leftmargin=1.4em]
\item An explicit framing of vendor changelog completeness as structurally unverifiable from within the system that authors it, together with a formal definition of a documented-vs-silent behavioral change as a coverage-conditioned lower bound rather than a claimed exact rate.
\item A calibrated diff detector that separates genuine cross-version behavioral change from ordinary sampling stochasticity via a corrected significance test and an effect-size floor, with its false-positive behavior checked directly against a same-version control.
\item A semantic-entailment changelog aligner --- validated against human labels and a naive string-matching floor --- that produces the first documented-or-silent label for a detected behavioral diff, rather than a bare change signal.
\item A behavior-category taxonomy anchored to independently published benchmarks, removing the circularity risk of a category scheme invented after seeing the data, together with a data-driven fingerprint check for attributing diffs to release windows that vendors do not cleanly date.
\item An empirical account, across multiple vendors and version transitions, of what fraction of measured behavioral change is silent and where it concentrates, reported honestly as a stability-checked lower bound.
\end{itemize}

\section{Problem Definition}

\noindent\textbf{Formal setup.} Consider a target system with consecutive versions $v_t, v_{t+1}$. A probe suite $P = \{p_1,\dots,p_M\}$ consists of fixed (prompt, decoding-configuration) pairs, each belonging to a category $c(p_i) \in \mathcal{C}$ (factual QA, refusal/safety, instruction-following/formatting, reasoning/math, tool-calling, latency/cost). For probe $p$ and version $v$, draw $N$ repeated samples $S_{p,v} = \{y_{p,v}^{(1)}, \dots, y_{p,v}^{(N)}\}$. The vendor's changelog for the transition is text $L_{t\to t+1}$, parsed into a discrete claim set $\mathcal{K} = \{k_1,\dots,k_R\}$. A probe is flagged as a \textbf{detected diff}, $D(p)=1$, if a within-version-null-corrected two-sample statistic on $S_{p,t}$ vs. $S_{p,t+1}$ exceeds significance after multiple-comparison correction and clears a minimum effect-size floor. A detected diff is \textbf{documented}, $\mathrm{Doc}(p)=1$, if some $k \in \mathcal{K}$ semantically entails it; it is \textbf{silent} otherwise: $\mathrm{Silent}(p) = D(p) \wedge \neg \mathrm{Doc}(p)$. The object of study is the resulting per-category and overall silent-change rate, $\rho_c$ and $\rho$, reported as a \textbf{coverage-conditioned lower bound} relative to the tested probe suite $P$ rather than an absolute rate over ``everything that actually changed.''

\noindent\textbf{Why this is hard.} Fixing $P$ does not fix a clean notion of $D(p)$: because same-version LLM sampling is itself variable (load-dependent routing, prompt sensitivity), the null distribution against which a cross-version difference in $S_{p,t}$ vs. $S_{p,t+1}$ is judged ``genuine'' must itself be estimated, not assumed, or an arbitrarily large probe count will manufacture significant-but-trivial diffs. Computing $\mathrm{Doc}(p)$ requires matching a probe's detected diff against $\mathcal{K}$, but $\mathcal{K}$ is extracted from unstructured, variable-granularity prose whose claims may entail many probes at once or none precisely --- a many-to-many semantic relationship, not a lookup. The category map $c(\cdot)$ that makes $\rho_c$ meaningful is only informative if $\mathcal{C}$'s boundaries were fixed independently of which probes turn out to satisfy $D(p)=1$; a category scheme built to explain an already-observed pattern would make any claim about $\rho_c$'s non-uniformity circular. And because $v_t, v_{t+1}$ are not always separated by a clean, dated release for API-served models, attributing $S_{p,v}$'s samples to the correct side of a specific transition --- so that $L_{t\to t+1}$ is even the right changelog to align against --- is itself uncertain rather than given. Finally, because no external party observes ``every change actually shipped'' in a release, $\rho$ and $\rho_c$ are only ever partially identified: they can be validated for \textbf{stability} under variation of $P$, but never certified as the true rate.

\section{Method Overview}

At a conceptual level, the study is a single measurement pipeline with five cooperating parts, each closing one of the difficulties above. Paired snapshots and a probe suite enter the pipeline; a calibrated detector turns the resulting samples into a genuine-diff flag, checked against a same-version control; a semantic aligner turns each flagged diff into a documented-or-silent label by comparison against the vendor's own changelog claims; a benchmark-anchored category scheme turns the resulting labels into a taxonomy that cannot be circular by construction; a fingerprint-based boundary check governs which transitions are clean enough to attribute diffs to in the first place; and a stability check turns the resulting rate from an unverifiable point claim into a validated lower bound.

The \textbf{calibrated diff detector} takes repeated samples from both snapshots of a probe, estimates the same-version null by resampling, and flags a genuine diff only when a corrected significance test and an effect-size floor are both cleared --- checked directly against a same-version control where no real diff should exist. The \textbf{entailment-based changelog aligner} takes each flagged diff and the parsed changelog claims for that transition and decides, by semantic entailment rather than string overlap, whether the diff is documented. The \textbf{category-anchored taxonomy} assigns each probe to a category defined by an independently published benchmark, so that a concentration pattern in the resulting per-category rates cannot be an artifact of how the categories were drawn. The \textbf{boundary fingerprinting} module continuously logs a lightweight distributional signature of the API and declares a transition ``clean'' only where that signature shows a genuine discontinuity, restricting headline claims to transitions it can confidently attribute. The \textbf{coverage stability check} recomputes the resulting rate from independent random subsets of the probe suite, reporting agreement rather than asserting the rate's absolute correctness.

\section{Methodology}

\subsection{Calibrated Diff Detector}
\textbf{Challenge.} Same-version LLM sampling is itself variable --- load-dependent routing, prompt sensitivity, and ordinary stochastic decoding mean that two samples drawn from the \emph{same} version already differ, so a fixed significance threshold alone is gameable by probe count and by trivially small effect sizes.

\textbf{Insight.} The same-version null is not a fixed constant to assume away but a distribution to estimate directly and check against; requiring both a corrected significance test and a minimum effect size, verified against a same-version control where no real change exists, turns ``is this a genuine diff'' into a falsifiable condition rather than an assumed one.

\textbf{Method.} For each probe, $N$ samples are drawn per snapshot; the same-version null is estimated by within-version resampling; a permutation/bootstrap two-sample test is applied at a fixed significance level, Holm-corrected across probes within a transition; a probe is only flagged $D(p)=1$ if it also clears a minimum standardized effect-size floor. The identical procedure is run on a same-version (A/A) resample pair as a null control, so the detector's own false-positive behavior is measured directly rather than assumed.

\subsection{Entailment-Based Changelog Alignment}
\textbf{Challenge.} Vendor changelogs are unstructured, variable-granularity prose --- a single claim can entail many probe categories at once, or none precisely --- so deciding whether a detected diff is covered by the changelog is an open semantic-matching problem, not a keyword lookup.

\textbf{Insight.} Only a semantic --- not lexical --- matcher can generalize across vendors' differing disclosure styles; benchmarking the semantic matcher against a naive string-matching floor is what makes ``semantic alignment earns its complexity'' a checkable claim instead of an assumed one.

\textbf{Method.} Changelog and release-note text for each transition is parsed into a discrete claim set. Each detected diff's description is scored for entailment against every claim in that set; a diff is labeled documented if its entailment score with any claim clears a threshold tuned on a human-labeled development subset, and silent otherwise. The same alignment is additionally computed with a naive keyword/string-overlap matcher, giving a lower-baseline comparison for how much of the labeling task is a lookup versus genuine inference.

\subsection{Category-Anchored Taxonomy}
\textbf{Challenge.} A behavior-category taxonomy invented after observing which probes showed silent changes would make any resulting ``concentration'' finding circular --- the categories could always be drawn to fit whatever pattern was already visible.

\textbf{Insight.} Anchoring each category to an independently published benchmark designed and validated for an unrelated original purpose breaks that circularity, because the category boundaries are fixed before, and independently of, any silent-change measurement.

\textbf{Method.} Each probe category is defined by a benchmark chosen for a purpose unrelated to this study: factual QA probes are drawn from an established knowledge-QA benchmark, reasoning/math probes from an established math-reasoning benchmark, hallucination probes from a truthfulness benchmark, over-refusal and under-refusal safety probes from complementary established safety benchmarks, formatting/instruction-following probes from an established instruction-compliance benchmark, and tool-calling probes from an established function-calling benchmark. Per-category silent-change rates are compared via a homogeneity test across categories, with pairwise comparisons Holm-corrected.

\subsection{Boundary Fingerprinting}
\textbf{Challenge.} API-served models are updated via rolling deploys, undated changes, and A/B routing rather than clean, discrete releases, so pairing ``before'' and ``after'' samples to the specific release window a changelog entry describes is itself uncertain.

\textbf{Insight.} An announced release date is not a reliable pairing device when deploys are rolling or undated; a discontinuity in a cheap, continuously logged behavioral fingerprint is a data-driven signal of a real boundary, checkable independently of what the vendor chooses to announce or when.

\textbf{Method.} The API is polled continuously and a lightweight distributional fingerprint (the first-token top-k log-probability profile) is logged over time. The fingerprint instrument itself is inherited from prior logprob-tracking work \citep{LogProbTracking} --- the first-token top-k log-probability profile is exactly the signal that line already validated for cheap change detection --- and no novelty is claimed for it; this paper's delta is its \emph{role}: the fingerprint's discontinuity serves as a transition-cleanliness gate deciding which transitions the disclosure-alignment claims may be attributed to at all. A transition is declared ``clean'' only where the fingerprint shows a discontinuity exceeding within-version fingerprint variance; headline silent-rate claims are restricted to clean transitions, with ambiguous or undated transitions quarantined and reported separately rather than silently pooled in. This module also fixes an eligibility constraint: it requires endpoints that expose token log-probabilities, which only a minority of production endpoints do, so vendor selection is screened for logprob availability before any transition is monitored.

\subsection{Coverage Stability Check}
\textbf{Challenge.} No independent record of ``everything that actually changed'' in a release exists, so the measured silent-change rate can never be certified as the true rate --- only ever as a lower bound conditioned on the tested probe suite's coverage.

\textbf{Insight.} Since the rate's absolute correctness is uncheckable, the closest available proxy for trustworthiness is stability: whether the same rate, or one close to it, is recovered from an independently sampled subset of the same probe suite.

\textbf{Method.} The probe suite for a given transition is split into two disjoint random halves, and the silent-change rate is recomputed independently on each half. Agreement between the two half-sample rates, within a pre-specified tolerance, is reported alongside the full-sample rate as the paper's stability evidence for that estimate, rather than as an implicit claim that the full-sample number is exactly correct.

\section{Experiments}

\subsection{Experimental Setup}

\textbf{Datasets and probe sources.} MMLU anchors the factual-QA category (an established, broad-coverage knowledge benchmark). GSM8K anchors reasoning/math. TruthfulQA anchors hallucination/factual-accuracy failure modes. XSTest anchors the over-refusal direction of safety behavior (safe prompts wrongly refused), while AdvBench/HarmBench anchor the complementary under-refusal direction (harmful prompts answered when they should be refused) --- both directions are needed because a vendor could silently shift refusal behavior either way between versions, and using only one source would leave the safety/refusal category blind to whichever direction it did not cover. IFEval anchors formatting/instruction-following compliance. BFCL anchors tool/function-calling behavior. A self-scraped changelog/release-note corpus spanning multiple API vendors is the documentation-of-record the pipeline aligns against.

\textbf{Snapshot acquisition and study-window gate.} Paired pre/post samples come from two sources: vendor-dated snapshot endpoints where they exist, and a prospective monitoring arm that polls each screened endpoint continuously so that both sides of an undated transition are captured as they occur --- retroactive measurement of a transition that predates monitoring is impossible, so the vendor set is fixed and monitoring started before any transition is scored. Because clean transitions arrive on the vendors' schedule rather than the study's, the design pre-registers a study-window gate: RQ5 is scored only once at least two fingerprint-clean and two quarantined transitions have accumulated across the vendor set; until then it is reported as pending rather than answered on an underpowered split.

\textbf{Baselines} (sanity-check comparators, not competitors to beat --- this is a measurement study, and no research question here is framed as ``our procedure beats baseline X''). A same-version (A/A) control, run through the identical pipeline on two same-version resamples, is the necessary check for every cross-version claim. A naive keyword/string-match changelog aligner is the weak lower baseline that justifies the semantic aligner's added complexity. Prior black-box change-detection tools (a logprob-distribution two-sample tracker; a rank-based model-substitution test) are run as change-detection-only comparators: if a probe's diff is flagged by neither prior detector, that is a sanity check on this pipeline's added coverage, not a claim of superiority over either.

\textbf{Metrics.} The detected-diff rate under cross-version and same-version (A/A) conditions; the overall and per-category silent-change rate ($\rho$, $\rho_c$) with bootstrap 95\% confidence intervals; entailment-aligner F1 against human labels and against the string-match floor; fingerprint-discontinuity status per transition; and split-half rate agreement.

\textbf{Research questions:}
\begin{enumerate}[leftmargin=1.6em]
\item RQ1 --- Does the diff detector control false positives against a same-version null while still detecting cross-version change?
\item RQ2 --- What fraction of detected behavioral diffs have no matching changelog claim?
\item RQ3 --- Does automatic changelog alignment reproduce human documented/silent labels, and does it improve on a naive string-match floor?
\item RQ4 --- Do silent changes concentrate in specific, externally-anchored behavior categories?
\item RQ5 --- Does restricting headline claims to fingerprint-clean version boundaries matter for the reported rate?
\item RQ6 --- Is the silent-change rate stable across independent random subsets of the probe suite?
\end{enumerate}

\noindent\textbf{Numbering crosswalk to the pre-registered protocol} (which numbers hypotheses H1--H3 rather than by RQ): RQ1 tests H3(b) (the A/A false-positive bound, gate PG-1); RQ2 tests H1 (the silent-rate majority claim, with its cross-vendor vote rule); RQ3 tests H3(a) (the alignment-F1 validity claim, gate PG-3); RQ4 tests H2 (the category-concentration claim). RQ5 and RQ6 carry no H-number: they are the pre-registered instrument-validity checks for boundary attribution (assumption A4, gate PG-4) and probe-coverage stability (assumption A5, gate PG-5), respectively.

\subsubsection{RQ1: Does the diff detector control false positives against a same-version null while still detecting cross-version change?}
\textbf{Setup.} Pilot stage: one vendor, one clean-boundary version transition, $\sim$200 probes across 5--6 categories, $N=50$ samples per snapshot; permutation two-sample test at $\alpha=0.05$, Holm-corrected across probes, with an effect-size floor of $|\text{Cliff's } \delta| \ge 0.33$. The identical procedure is additionally run on a same-version (A/A) resample pair at equal sampling budget. Scale-up: $\geq 3$ vendors $\times$ multiple transitions each.

\textbf{Expected outcome.} The A/A control flags detected diffs ($D=1$) at a rate at or below the nominal $\alpha=0.05$; the cross-version arm's count of detected diffs exceeds the A/A arm's false-diff count by at least $3\times$.

\textbf{Falsification condition.} Refuted if the A/A control's false-diff rate exceeds $\alpha$ (the detector manufactures its own findings even absent real change), or if the cross-version arm's detected-diff count does not exceed the A/A arm's by the pre-registered $3\times$ margin (cross-version ``diffs'' would then be statistically indistinguishable from the detector's own noise floor) --- either would mean the headline silent-change rate is an artifact of the detector rather than of real vendor behavior. Per the pre-registered fallback, an A/A false-diff rate exceeding $\alpha$ does not end the study: the effect-size floor is recalibrated upward to the smallest value that restores FPR $\le \alpha$ on the A/A arm, all rates are re-derived under the corrected floor, and the corrected floor is reported transparently alongside the original --- never silently substituted.

\textbf{Hypothesis (plain language).} The pipeline can tell a genuine cross-version behavioral change apart from ordinary same-version sampling noise, not just claim to.

\textbf{Metrics.} Detected-diff rate and count, cross-version vs. A/A; margin ratio between the two arms.

\textbf{Result.} \ph

\subsubsection{RQ2: What fraction of detected behavioral diffs have no matching changelog claim?}
\textbf{Setup.} For every probe flagged $D(p)=1$ in a transition, compute the documented/silent label via the entailment aligner against that transition's parsed changelog claims; aggregate into an overall rate $\rho$ with bootstrap 95\% confidence intervals ($\geq 1000$ resamples), reported per vendor rather than pooled. Because $\rho$ is scored per vendor rather than pooled, H1 is judged by a pre-registered \textbf{cross-vendor vote}: PASS requires $\rho>0.50$ with CI lower bound also $>0.50$ in at least 2 of the $\geq3$ tested vendors (a vendor contributing fewer than 20 detected diffs is excluded from the vote as underpowered, not counted as a ``no''). At Stage-0 pilot scale, with only one vendor tested, H1 is not scored PASS/FAIL at all --- it is reported \textbf{INCONCLUSIVE-pending-scale-up} until the $\geq3$-vendor vote is available.

\textbf{Expected outcome.} The cross-vendor vote passes: $\rho$ exceeds 0.50 with the bootstrap 95\% CI lower bound also exceeding 0.50 in at least 2 of the $\geq3$ tested vendors.

\textbf{Falsification condition.} Refuted if the cross-vendor vote fails (fewer than 2 of the $\geq3$ tested vendors clear $\rho>0.50$ with CI lower bound $>0.50$) --- in that case the central premise (that silence is the majority outcome) does not hold across the tested vendor set, and the paper's claim shifts, per the pre-committed fallback, from ``silence is the default'' to ``silence is a material minority,'' reporting the measured $\rho$ and its CI directly, per vendor, rather than the majority framing, with the headline reframed around RQ4's existence-and-category-structure finding rather than majority prevalence.

\textbf{Hypothesis (plain language).} For the systems tested, most statistically-detected behavioral change is undocumented, not the exception.

\textbf{Metrics.} Overall silent-change rate $\rho$ with bootstrap CI, per vendor.

\textbf{Result.} \ph

\subsubsection{RQ3: Does automatic changelog alignment reproduce human documented/silent labels, and does it improve on a naive string-match floor?}
\textbf{Setup.} A human-labeled subset of $\geq 200$ diff/claim pairs is adjudicated by two annotators with disagreement resolution, reporting Cohen's $\kappa$ (required $\geq 0.6$ for label validity). The entailment aligner's documented/silent labels are compared against this human-labeled subset via F1, and against the naive string-match aligner's F1 on the same subset.

\textbf{Expected outcome.} The entailment aligner achieves $F1 \geq 0.70$ against human labels, exceeding the string-match baseline's F1 by a pre-specified, nontrivial margin.

\textbf{Falsification condition.} Refuted (alignment is unreliable) if entailment-aligner F1 falls below 0.70; refuted (semantic matching is unnecessary complexity) if the string-match baseline's F1 is statistically indistinguishable from the entailment aligner's --- either would undercut the claim that this paper's alignment step is doing more than a lookup. Per the pre-registered fallback, an F1 below 0.70 does not retract M4: the automatic aligner is downgraded from a headline instrument to a human-in-the-loop aid (auto-flag plus mandatory human adjudication), and the silent-change rate is reported only on the human-labeled subset, with the automatic-pipeline claim explicitly relabeled ``assistive'' rather than validated.

\textbf{Hypothesis (plain language).} Deciding whether a changelog covers a detected behavioral change is a genuine inference problem that a keyword search cannot substitute for.

\textbf{Metrics.} F1 (entailment vs. human, string-match vs. human), Cohen's $\kappa$ on the human-labeled subset.

\textbf{Result.} \ph

\subsubsection{RQ4: Do silent changes concentrate in specific, externally-anchored behavior categories?}
\textbf{Setup.} Per-category silent-change rate $\rho_c$ is computed for each benchmark-anchored category with $\geq 20$ detected diffs (categories below this threshold are pooled and reported as underpowered rather than tested individually). A chi-square homogeneity test checks whether $\rho_c$ is uniform across categories; Holm-corrected pairwise proportion tests check the specific predicted ordering (safety/refusal, formatting/instruction-following, and latency/cost each above factual-accuracy).

\textbf{Expected outcome.} $\rho_c$ is non-uniform across categories (chi-square rejects equality at $\alpha=0.05$), with safety/refusal, formatting/instruction-following, and latency/cost each significantly exceeding factual-accuracy in the pairwise tests.

\textbf{Falsification condition.} Refuted (pervasive, category-agnostic silence) if the homogeneity test does not reject uniformity; refuted in direction (a different concentration pattern) if $\rho_c$ is non-uniform but the predicted ordering does not hold --- in the latter case the observed concentration, whatever it is, is reported as the finding, since the taxonomy's value is in surfacing genuine structure, not in confirming a specific predicted order.

\textbf{Hypothesis (plain language).} Vendors are not equally likely to leave every kind of behavior change undocumented --- some behavior categories are silently changed far more often than others.

\textbf{Metrics.} Per-category $\rho_c$ with bootstrap CIs; chi-square statistic; pairwise proportion test results.

\textbf{Result.} \ph

\subsubsection{RQ5: Does restricting headline claims to fingerprint-clean version boundaries matter for the reported rate?}
\textbf{Setup.} Continuous fingerprint logging classifies each candidate transition as clean (fingerprint discontinuity exceeds within-version variance) or quarantined (no clear discontinuity, e.g., undated or rolling changes). The detected-diff rate is compared between clean and quarantined transitions, and the headline silent-change rate computed on clean transitions only is compared against a naive estimate pooling all transitions including quarantined ones.

\textbf{Expected outcome.} Clean transitions show a significantly higher detected-diff rate than quarantined ones, and the clean-only headline rate differs materially (not just with wider uncertainty) from the naive pooled estimate.

\textbf{Falsification condition.} Refuted if detected-diff rates are statistically indistinguishable between clean and quarantined transitions, or if the clean-only and naive pooled rates agree within noise --- either would mean the fingerprint-based boundary check is not doing load-bearing work and a simpler, calendar-based pairing would have sufficed.

\textbf{Hypothesis (plain language).} Attributing diffs to the release window a changelog actually describes requires a data-driven boundary check, not just trusting the vendor's announced release date.

\textbf{Metrics.} Detected-diff rate by transition-cleanliness status; clean-only vs. pooled silent-change rate comparison.

\textbf{Result.} \ph

\subsubsection{RQ6: Is the silent-change rate stable across independent random subsets of the probe suite?}
\textbf{Setup.} For each (vendor $\times$ transition) cell with sufficient detected diffs, the probe suite is split into two disjoint random halves and the silent-change rate recomputed independently on each half; agreement between the two half-sample rates is compared against a pre-specified tolerance.

\textbf{Expected outcome.} The two half-sample rates agree within the pre-specified tolerance in a majority of tested cells.

\textbf{Falsification condition.} Refuted if agreement falls below the majority threshold across tested cells --- that would mean the reported rate is sensitive to which probes happened to be included, undermining its interpretation as a property of the vendor's disclosure practice rather than of an arbitrarily chosen probe list.

\textbf{Hypothesis (plain language).} The reported silent-change rate describes something real about how a vendor documents change, not an artifact of which specific probes were chosen to measure it.

\textbf{Metrics.} Split-half rate agreement rate across tested cells.

\textbf{Result.} \ph

\vspace{0.5em}
\noindent\textbf{Further evaluation.} With more time, a natural extension is adversarial robustness: testing whether a vendor aware of being audited could serve an audit-facing response distribution distinct from its user-facing one, undermining the entire measurement from outside the pipeline described here. This is explicitly out of scope for the present design, which assumes the auditor's queries are indistinguishable from ordinary traffic; detecting or defending against audit-aware serving would require its own identifiability design, not an extension of the modules above. Recent work on fingerprint spoofing \citep[arXiv:2606.16100]{Zhang2026FingerprintSpoof} makes this limitation concrete: the same lightweight logprob fingerprints this design relies on for boundary attribution can in principle be imitated by a serving stack that wishes to hide a transition --- so the boundary-fingerprinting module, like every module here, is sound only under the stated non-adversarial-serving assumption.

\section{Related Work}

The closest prior work to this paper is Chen, Zaharia, and Zou's study of ChatGPT's behavioral drift over time \citep{Chen2023ChatGPTBehavior}, and it is worth stating the delta plainly before surveying the wider landscape: that work already demonstrates substantial, measurable behavioral drift between two dated GPT-3.5/GPT-4 snapshots on a fixed task battery, but it never cross-references the vendor's own release notes to determine which of the observed changes were disclosed and which were silent --- precisely the alignment step this paper is built to supply, and the reason a documented-vs-silent rate has never been reported for any deployed system, including the one that closest work already measured.

\textbf{Black-box API change and model-equality detection.} Continuous logprob-tracking methods \citep{LogProbTracking} repeatedly query an endpoint with a fixed minimal prompt and apply a nonparametric two-sample test to the resulting top-k token log-probabilities, giving a cheap, general signal that a served distribution has shifted. That line's field deployment is the nearest published evidence to this paper's headline question: monitoring 189 production endpoints across ten providers for over four months, it reports 37 suspected undisclosed changes --- but ``undisclosed'' there is a presumption inherited from version pinning (a pinned endpoint should never change at all), not a determination made by parsing the vendor's disclosure text, which that work itself names as outside its scope; and a token-efficient successor in the same line \citep[arXiv:2602.11083]{Chauvin2026TokenEfficient} reduces the query cost of the same detection signal without adding any disclosure arm. Rank-based uniformity tests \citep{RankUniformity} instead ask a narrower, binary question --- is the served model identical to a trusted local reference --- by scoring completions' percentile ranks under that reference and testing for uniformity, flagging silent model substitution. A related line traces a fine-tuned model's black-box outputs back to its pre-trained base \citep{Foley2023Matching}, the same family of concern (an undocumented substitution behind an API) solved via lineage attribution rather than disclosure alignment. All of these establish that behavioral change or substitution in a deployed API is detectable at all, but none attributes a detected shift to a behavior category or aligns it against what the vendor actually disclosed; a bare change signal --- even one presumed undisclosed --- is not a documented-or-silent label, and no documented-vs-silent rate has been reported.

\textbf{Relational and fairness-specific behavior audits.} One line of work measures whether a fine-tuned release departs from its base model by representing each model's behavior as a multi-view signature over shared fairness prompts and computing pairwise similarity \citep{BiasSimilarity}, detecting inheritance and drift within the bias/stereotype construct specifically. A related line shows that a single-time audit under this kind of construct is vulnerable to a platform serving a compliant audit-facing model distinct from its user-facing one, and proposes using an auditor's private labeled prior to make such manipulation statistically detectable \citep{RobustAudit}. Both are confined to a fairness-specific construct and a one-shot audit design, neither spans the general behavior categories this paper's taxonomy anchors nor tracks documentation fidelity across a longitudinal sequence of releases.

\textbf{Governance positions and accountability infrastructure.} A survey of technical AI governance \citep{TAIG} decomposes regulatory mandates --- registration, evaluation, auditing, and especially ongoing monitoring --- into concrete technical requirements and catalogs where current tooling falls short of satisfying them; a companion position paper \citep{SafeHarbor} argues that independent third-party evaluation is a structural prerequisite for any of this to happen, given that Terms-of-Service restrictions make such evaluation legally hostile today, and proposes paired legal and technical safe harbors. An organizational-risk framework for AI-native engineering teams \citep{RiskArchitecture2607} independently names the same accountability gap from the deployer's side --- "providers update weights, third-party tool contracts drift... with no corresponding change event" --- and a standardized-protocol proposal for LLM-judge/evaluator drift \citep{EPC2607} raises the identical structural concern one layer up the stack, for evaluators rather than the systems they evaluate, solving it via community versioning rather than changelog-text alignment. None of these four builds or runs an empirical instrument --- they motivate the audit without operationalizing it.

\textbf{Longitudinal drift measurement and sequential statistical machinery.} A martingale-based betting framework \citep{FairBetting} gives anytime-valid sequential tests for a group-fairness null that remain correct under continuous monitoring and optional stopping --- reusable statistical machinery for a longitudinal audit, but with no mechanism to parse prose changelogs or build a behavior-category taxonomy. A decade-spanning malware benchmark \citep{LAMDA} demonstrates rigorous longitudinal concept-drift evaluation under proper temporal splits, but for a first-party classifier's own accuracy degradation, not the gap between a black-box vendor's behavior and that vendor's own documentation. Together with the direct precedent discussed above, these establish that every statistical and methodological ingredient this paper needs exists somewhere in the literature --- but none of them has been pointed at the vendor's own disclosure text.

\section{Conclusion}

This work reframes vendor changelog completeness --- usually taken for granted as the record of what changed between API versions --- as an empirically measurable quantity rather than an assumption, because the party who controls when a system changes is structurally the same party who authors the record of that change. The central methodological commitment is that this measurement is only credible if the instrument's own properties are checked rather than assumed: whether the diff detector can be fooled by same-version noise, whether the changelog aligner does more than a keyword search, whether the category taxonomy could have manufactured its own concentration pattern, whether version-boundary attribution matters, and whether the headline rate is stable rather than an artifact of a particular probe list. Whichever silent-change rate is measured, the finding is decision-relevant either way: a high rate means changelog-based accountability is currently unreliable and needs an external, recurring check exactly like this one; a low rate would itself be a notable finding about the tested vendors' disclosure practices, and either result is reported as a coverage-conditioned lower bound rather than a precise absolute figure. Two limitations bound this claim from the outset. First, the measured rate can never be certified as the true rate, because no independent record of ``everything that actually changed'' in a release will ever exist --- this design reports and stability-checks a lower bound rather than overclaiming precision it cannot have. Second, findings from the tested vendors may not generalize to vendors with different documentation cultures, release cadences, or serving infrastructure --- and the tested set is itself screened for logprob availability, so it is not a random sample of the market --- which is why per-vendor rates are reported separately rather than pooled into a single number; extending the audited vendor set is the natural next step, not an afterthought. If the central hypothesis family fails outright --- the cross-vendor vote for RQ2 does not pass and RQ4's category rates turn out uniform --- the contribution does not disappear: it pivots, per a pre-committed plan, from a measurement-finding paper to an instrument paper, with the validated, error-characterized documentation-fidelity pipeline itself as the headline contribution and the null result reported as an accountability-relevant calibration finding (vendor changelogs for the tested systems are largely faithful) rather than a failed study.

\vspace{1em}
\hrule
\vspace{1em}

\section*{Pre-handoff checklist}
\begin{itemize}[leftmargin=1.4em]
\item \textbf{Organization/navigation} --- narrative flows problem $\to$ method $\to$ experiments $\to$ related work without forward references to undefined terms.
\item \textbf{Terminology discipline} --- ``detected diff,'' ``documented/silent,'' ``silent-change rate,'' ``coverage-conditioned lower bound'' each defined at first use in Problem Definition, used consistently thereafter; no internal stage IDs (\texttt{sp\_*}, \texttt{m\_*}, \texttt{rq\_*}) leaked into prose --- checked directly against the discussion pipeline log.
\item \textbf{Notation with worked example} --- the formal setup section defines $(v_t, v_{t+1}, P, c(\cdot), S_{p,v}, L, \mathcal{K}, D(p), \mathrm{Doc}(p), \mathrm{Silent}(p), \rho, \rho_c)$ and the Problem Definition/Methodology sections consistently reuse this notation; no separate worked numeric example is included since no numbers exist yet (would require fabricating results).
\item \textbf{RQ framing matches outcomes} --- all 6 RQs are graded/quantitative (rates, F1, agreement proportions), each answered in its own subsection with an explicit falsifiable expected outcome and a stated falsification condition, not a binary yes/no.
\item \textbf{No typesetting artifacts} --- only explicit placeholder markers in Results, per the drafting convention this is a pre-registered spec, not a report.
\item \textbf{Artifact plan} --- the probe suite, the version-fingerprinting logger, the entailment-alignment code, and the frozen raw-response audit trail are the natural release artifacts; not yet built, so not claimed as available.
\end{itemize}

\section*{Hand-off}
This is a draft, not a finished paper --- experiments have not run. Natural next steps: prose polish and venue-specific LaTeX formatting once a target venue is chosen, or revising an earlier stage if piloting surfaces a problem (e.g., if the string-match baseline in RQ3 is closer to the entailment aligner's F1 than expected).

\bibliographystyle{plainnat}
\bibliography{refs}

\end{document}
